\documentclass[11pt]{scrartcl}
\usepackage{evan}
\ihead{Berkeley Math Circle}
\ohead{Parallelograms}

\usepackage{wrapfig}
\begin{asydef}
  // Old-style angle mark and tick-mark helpers, matching the long-retired
  // olympiad.asy versions.
  real markscalefactor = 0.03;
  path anglemark(pair A, pair B, pair C, real t = 8) {
    pair M = t * markscalefactor * unit(A - B) + B;
    pair N = t * markscalefactor * unit(C - B) + B;
    return arc(B, M, N) -- B -- cycle;
  }
  picture pathticks(path g, int n = 1, real r = .5, real spacing = 6,
      real s = 8, pen p = currentpen) {
    picture pict;
    real l = arclength(g);
    real space = spacing * markscalefactor;
    real halftick = s * markscalefactor / 2;
    if (n > 0) {
      pair direct = unit(dir(g, arctime(g, r * l)));
      real startpt = r * l - (n - 1) / 2 * space;
      for (int i = 0; i < n; ++i) {
        real t = startpt + i * space;
        pair B = point(g, arctime(g, t)) + (0, 1) * halftick * direct;
        pair C = B + 2 * (0, -1) * halftick * direct;
        draw(pict, B -- C, p);
      }
    }
    return pict;
  }

  real m = 1;
  pair bisect(pair X, pair Y, int n=2, real r=0.5, real spacing=4*m, real s=7*m, pen p=currentpen) {
    pair M = midpoint(X--Y);
    add(pathticks(X--M, n, r, spacing, s, p));
    add(pathticks(Y--M, n, r, spacing, s, p));
    return M;
  }
\end{asydef}

\renewcommand{\figurename}{Problem}
\renewcommand{\thefigure}{\thetheorem}

\begin{document}
\title{All you have to do is construct a parallelogram!}
\author{Evan Chen\plusemail{evanchen@mit.edu}}
\date{Berkeley Math Circle \\ January 9, 2014}
\maketitle

\begin{abstract}
  As the name suggests, all of these problems can be solved by constructing the fourth vertex of a parallelogram. Most require just one point to be drawn in. A few will require constructing multiple points or even multiple parallelograms.
\end{abstract}


\section{Samples}
\begin{problem}
  Let $M$ and $N$ be the midpoints of $\ol{AB}$ and $\ol{AC}$ in triangle $ABC$.  Prove $MN = \half BC$ without using similar triangles.
\end{problem}



\section{Appetizers}
\begin{problem}
  Let $M$ be the midpoint of $\ol{BC}$ in a triangle $ABC$.  Given that $AM=2$, $AB = 3$, $AC = 4$, find the area of $ABC$.
\end{problem}

\begin{problem}[AIME 2011]
  In rectangle $ABCD$, $AB=12$ and $BC=10$. Points $E$ and $F$ lie inside rectangle $ABCD$ so that $BE=9$, $DF=8$, $\overline{BE} \parallel \overline{DF}$, $\overline{EF} \parallel \overline{AB}$, and line $BE$ intersects segment $\overline{AD}$.  Find $EF$.
\end{problem}

\begin{problem}
  Let $ABC$ be a triangle and $M$ be the midpoint of $\ol{BC}$.  Squares $ABQP$ and $ACYX$ are erected. Show that $PX = 2AM$.
\end{problem}

\begin{figure}[ht]
  \centering
  \begin{minipage}{6cm}
  \centering
  \begin{asy}
    size(4cm);
    pair D = Drawing("D", (5,7), dir(45));
    pair A = Drawing("A", (0,D.y), dir(135));
    pair B = Drawing("B", origin, dir(225));
    pair C = Drawing("C", (D.x,0), dir(-45));
    pair E = Drawing("E", (3,6), dir(180));
    pair F = Drawing("F", (3,1), dir(0));
    draw(A--B--C--D--cycle);
    draw(B--E--F--D);
  \end{asy}
  \addtocounter{theorem}{-1}
  \caption{AIME 2011.}
  \stepcounter{theorem}
  \end{minipage}
  \begin{minipage}{8cm}
  \begin{asy}
    size(7cm);
    pair A = Drawing("A", origin, dir(110));
    pair B = Drawing("B", (-0.8,-1.9), dir(-90));
    pair C = Drawing("C", (1.8,-1.9), dir(-90));
    pair P = Drawing("P", rotate(-90)*B, dir(160));
    pair Q = Drawing("Q", P+B, dir(P+B));
    pair X = Drawing("X", rotate(90)*C, dir(60));
    pair Y = Drawing("Y", X+C, dir(X+C));
    pair M = Drawing("M", bisect(B,C), dir(-90));
    draw(B--Q--P--A--X--Y--C--B--A--C);
    draw(P--X, dotted); draw(A--M, dotted);
  \end{asy}
  \caption{If $APQB$ and $AXYC$ are squares, prove $PX=2AM$.}
\end{minipage}
\end{figure}



\section{Meals}
\begin{problem}
  The area of triangle $ABC$ is $4$, and the length of the medians are $m_a$, $m_b$, and $m_c$. A second triangle has side lengths $m_a$, $m_b$, and $m_c$.  What is its area?
\end{problem}

\begin{problem}[USAMO 2003]
  Let $ABC$ be a triangle. A circle passing through $A$ and $B$ intersects segments $AC$ and $BC$ at $D$ and $E$, respectively. Lines $AB$ and $DE$ intersect at $F$, while lines $BD$ and $CF$ intersect at $M$. Prove that $MF = MC$ if and only if $MB\cdot MD = MC^2$.
\end{problem}

\begin{problem}[NIMO 8.8]
  The diagonals of convex quadrilateral $BSCT$ meet at the midpoint $M$ of $\overline{ST}$.  Lines $BT$ and $SC$ meet at $A$, and $AB = 91$, $BC = 98$, $CA = 105$.  Given that $\overline{AM} \perp \overline{BC}$, find the positive difference between the areas of $\triangle SMC$ and $\triangle BMT$.
\end{problem}

\begin{figure}[ht]
  \centering
  \begin{minipage}{8cm}
  \begin{asy}
    size(7.5cm);
    pair C = Drawing("C", (3,0), dir(-45));
    pair D = Drawing("D", (0,0));
    pair M = Drawing("M", (1,-0.7));
    pair F = Drawing("F", 2*M-C);
    pair B = Drawing("B", conj(C/M) * (M-C)+C, dir(130));
    pair E = Drawing("E", extension(D,F,B,C), dir(40));
    pair A = Drawing("A", extension(D,C,B,F), dir(225));
    draw(A--B--C--cycle);
    draw(circumcircle(B,D,E), dashed);
    draw(A--F, dotted);
    draw(E--F--C, dotted);
    draw(B--M, dotted);
    bisect(F,C);
  \end{asy}
  \addtocounter{theorem}{-1}
  \caption{USAMO 2003.}
  \stepcounter{theorem}
  \end{minipage}
  \begin{minipage}{6cm}
  \begin{asy}
    size(5.5cm);
    pair M = Drawing("M", (0,0), dir(-90));
    pair S = Drawing("S", (2,1.6), dir(45));
    pair T = Drawing("T", -S, dir(225));
    pair A = Drawing("A", (0,3.7), dir(90));
    pair z = (1,0);
    pair B = Drawing("B", extension(M,z,A,T), dir(180));
    pair C = Drawing("C", extension(M,z,A,S), dir(0));
    draw(A--B--C--cycle);
    bisect(S,T);
    draw(A--M);
    draw(B--T, dotted);
    draw(S--T);
    draw(B--S, dashed);
    draw(T--C, dashed);
  \end{asy}
  \caption{NIMO 8.8.}
  \end{minipage}
\end{figure}

\begin{problem}
  Let $ABC$ be a triangle with circumcenter $O$ and orthocenter $H$, and let $M$ and $N$ be the midpoints of $\ol{AB}$ and $\ol{AC}$.  Rays $MO$ and $NO$ meet line $BC$ at $Y$ and $X$, respectively.  Lines $MX$ and $NY$ meet at $P$.  Prove that $\ol{OP}$ bisects $\ol{AH}$.
\end{problem}

\begin{figure}[h]
  \centering
  \begin{asy}
    size(7.5cm);
    pair A = Drawing("A", dir(110), dir(110));
    pair B = Drawing("B", dir(210), dir(-90));
    pair C = Drawing("C", dir(330), dir(-90));
    pair H = Drawing("H", orthocenter(A,B,C), dir(135));
    pair O = Drawing("O", circumcenter(A,B,C), dir(-90));
    pair M = Drawing("M", midpoint(A--B), dir(A+B));
    pair N = Drawing("N", midpoint(A--C), dir(A+C));
    pair X = Drawing("X", extension(O,N,B,C), dir(-90));
    pair Y = Drawing("Y", extension(O,M,B,C), dir(-90));
    pair P = Drawing("P", extension(M,X,N,Y), dir(180));
    draw(A--B--C--cycle);
    draw(A--H);
    draw(P--O);
    draw(M--Y, dashed);
    draw(N--X, dashed);
    draw(X--P--Y, dotted);
    draw(C--Y, dotted);
    markscalefactor *= 0.3;
    dot(bisect(A,H));
  \end{asy}
  \caption{Show that $\ol{PO}$ bisects $\ol{AH}$.}
\end{figure}



\section{Buffets}
\begin{problem}
  Let $ABC$ be a triangle with orthocenter $H$ and let $P$ be a point on the circumcircle of $ABC$.  The \emph{Simson line} from $P$ is the line passing through the feet of the altitudes from $P$ to $\ol{BC}$, $\ol{CA}$, $\ol{AB}$.  Prove that it bisects $\ol{PH}$.
\end{problem}
You may want to use this lemma: the reflection of $H$ over $\ol{BC}$ lies on the circumcircle.

\begin{figure}[ht]
  \centering
  \begin{asy}
    size(7cm);
    pair A = dir(110);
    pair B = dir(200);
    pair C = dir(340);
    pair P = dir(50);
    Drawing("A", A, A);
    Drawing("B", B, B);
    Drawing("C", C, C);
    Drawing("P", P, P);
    pair X = Drawing(foot(P,B,C));
    pair Y = Drawing(foot(P,C,A));
    pair Z = Drawing(foot(P,A,B));
    draw(P--X, dotted);
    draw(P--Y, dotted);
    draw(P--Z, dotted);
    draw(A--B--C--cycle);
    pair H = Drawing("H", orthocenter(A,B,C), dir(225));
    draw(Line(X,Z,0.3,0.2), black, Arrows);
    draw(P--H);
    draw(A--Z, dashed);
    draw(unitcircle, dashed);
    markscalefactor*=0.3;
    dot(bisect(P,H));
  \end{asy}
  \caption{The Simson line bisects $\ol{PH}$.}
\end{figure}

\begin{problem}
  Let $ABCDE$ be a convex pentagon with $AB=BC$ and $CD=DE$.  If $\angle ABC = 2\angle CDE = 120\dg$ and $BD=2$, find the area of $ABCDE$.
\end{problem}
\begin{figure}[ht]
  \centering
  \begin{asy}
    size(6cm);
    pair D = Drawing("D", (0,0), dir(180));
    pair B = Drawing("B", (1,0), dir(0));
    pair C = Drawing("C", (0.6,-0.2), dir(-90));
    pair E = Drawing("E", rotate(60) * C, dir(90));
    pair A = Drawing("A", rotate(-120,B)*C, dir(90));
    draw(A--B--C--D--E--cycle);
    draw(B--D);
    MP("2", foot(C,B,D), dir(90));
    add(pathticks(C--D, 1, 0.5, 1, 1.5));
    add(pathticks(E--D, 1, 0.5, 1, 1.5));
    add(pathticks(A--B, 2, 0.5, 1, 1.5));
    add(pathticks(C--B, 2, 0.5, 1, 1.5));
    markscalefactor *= 0.4;
    MP("60^{\circ}", D, dir(19)*5);
    draw(anglemark(C,D,E));
    MP("120^{\circ}", B, dir(130)*4);
    draw(anglemark(A,B,C));
  \end{asy}
  \caption{If $\angle ABC = 2\angle CDE = 120\dg$ and $BD=2$, find the area of $ABCDE$.}
\end{figure}

\begin{problem}
  [ELMO 2012] Let $ABC$ be an acute triangle with $AB<AC$, and let $D$ and $E$ be points on side $BC$ such that $BD=CE$ and $D$ lies between $B$ and $E$. Suppose there exists a point $P$ inside $ABC$ such that $\ol{PD} \parallel \ol{AE}$ and $\angle PAB=\angle EAC$.
  Prove that $\angle PBA = \angle PCA$.
\end{problem}



\end{document}
